% Chapter: Introduction
% Porous BZ reactor as a structured reaction-diffusion computing substrate

\section{Background and Motivation}
\label{sec:background}

Digital computers move data between memory and processor for every operation, and this von Neumann bottleneck costs energy and latency at every arithmetic step.
As machine-learning workloads grow, there is increasing interest in using the physics of a material to compute directly, without clocked digital arithmetic.
In such physical neural networks, inputs are encoded as physical excitations, the medium transforms them through its own dynamics, and outputs are read from probes.
Computation is continuous in time and parallel across space, with the medium's structure acting as the network's weights.
Demonstrations in optics, mechanics, and electronics show that this approach can work in practice \cite{lin2018alloptical, wright2022deep, hughes2019wave}.

This project starts from a simple question: what happens inside a porous chemical reactor when a pulse of reagent is injected at one face?
The reactor is a solid scaffold saturated with a reacting fluid; its internal geometry is a maze of pores joined by narrow throats.
A chemical excitation introduced at one pore will travel through the fluid, split at junctions, collide with other excitations, and possibly reach output locations on another face.
If the local rules that govern this traffic are known, the reactor can be treated as a programmable analog computer whose program is written into the pore geometry and the spatial distribution of chemical properties.

Making this precise requires two things: a medium rich enough to support directed pulse traffic, and a systematic way to describe what a given arrangement of it computes.
This thesis builds the case for one answer and tests it.
\refSec{sec:lit_pnn} reviews the physical-neural-network idea, \refSec{sec:lit_chemical} the chemical and excitable-media route to computation, and \refSec{sec:lit_bz} the Belousov--Zhabotinsky reaction and Oregonator model as the practical realisation.
\refSec{sec:lit_synthesis} draws these threads into the research gap, from which the aim and objectives follow.

\section{Physical Neural Networks}
\label{sec:lit_pnn}

A physical neural network is a material or device whose internal physical dynamics implement the operations of a neural network.
Inputs are encoded as fields or excitations, the material transforms them through its natural evolution, and outputs are read from detectors.
The appeal is that computation can occur in parallel across space and continuously in time, without the repeated memory transfers that dominate digital inference.

Optical systems provide the most prominent demonstrations.
Silva \etal showed that a metamaterial slab can perform mathematical operations such as spatial differentiation and integration on an incident electromagnetic field \cite{silva2014performing}.
Shen \etal demonstrated a coherent nanophotonic circuit that implements a deep neural network for image classification entirely in the optical domain \cite{shen2017deep}.
Lin \etal extended this to diffractive deep neural networks, in which successive phase masks modulate free-space propagation so that the outgoing intensity pattern encodes the classification result \cite{lin2018alloptical}.
The geometry of the masks or scatterers serves as the trainable weights.

Mechanical and thermal analogs offer complementary physical bases.
Stern \etal demonstrated supervised learning through physical changes in a mechanical network of springs and masses \cite{stern2020supervised}, and later developed a broader framework for learning in physical networks \cite{stern2021physical}.
Wright \etal extended the idea to deep architectures by chaining multiple physical systems and training them with backpropagation through differentiable simulations \cite{wright2022deep}.
Thermal analog computing has been proposed for matrix-vector multiplication using inverse-designed metastructures \cite{silva2026thermal}.
These demonstrations prove that the PNN concept is not restricted to any one physics.

A unifying theoretical perspective was provided by Hughes \etal, who showed that linear wave propagation can be written as an analog recurrent neural network \cite{hughes2019wave}.
In this view, the wave field at each time step is the hidden state, and the evolution operator is a sparse linear recurrence whose coefficients depend on local material properties.
McMahon reviewed the broader trend of nonlinear computation with linear systems and argued that such approaches are becoming practical for specialised tasks \cite{mcmahon2024nonlinear}.

The present work adopts the same stance---the physical field is the computational state and the material distribution is the trainable parameter set---but moves from linear wave physics to a nonlinear, excitable chemical medium.
The propagating quantity is a self-sustaining chemical pulse rather than a decaying or superposing wave.
This situates the work within molecular communication: the substrate is a fluid channel whose input is a spatially patterned chemical excitation and whose output is a measured concentration signal.
Experimental tabletop molecular communication has already demonstrated digital messages carried by chemical signals \cite{farsad2013tabletop}, and surveys identify reaction--diffusion transport as a physically realistic channel model \cite{farsad2016comprehensive}.
The present thesis treats the reaction--diffusion channel not merely as a passive link, but as a programmable analog computing substrate.

\section{Chemical and Excitable-Media Computing}
\label{sec:lit_chemical}

Computation in fluid and chemical media has a long history.
Matter itself stores and transforms information, so signals are processed by the physics of the substrate rather than by an external arithmetic unit.
The lineage runs from hydraulic analog machines to modern molecular and microfluidic systems \cite{adamatzky2019liquid}.
Demonstrated platforms include microfluidic bubble logic, in which gas bubbles carry, route, and process discrete bits through channel networks \cite{prakash2007bubble}, and self-organising chemical reaction networks that classify temporal inputs as reservoir computers \cite{baltussen2024chemical}.

Within this tradition, excitable chemical media offer a direct route to computation because their elementary phenomena map onto computational primitives.
An excitable medium has a stable rest state, a threshold for excitation, and a refractory period after excitation during which the medium cannot be re-excited.
A supra-threshold perturbation triggers a self-sustaining pulse that propagates at a characteristic speed, leaving a refractory wake.
Colliding pulses annihilate; pulses meeting a narrowing gap or high-curvature turn can be blocked \cite{courtemanche1993instabilities}.
These primitives have produced working devices: colliding-pulse logic gates, an experimental XOR gate in a reaction--diffusion medium, logical functions of structured channel junctions, and a single oscillating BZ droplet acting as a binary and fuzzy logic processor \cite{toth1995logic, steinbock1996chemical, adamatzky2002xor, adamatzky2002collision, sielewiesiuk2001logical, gorecki2009information, gentili2012chemical}.
The common feature is hand design: each gate or circuit is built geometry by geometry for one logical function.

A parallel line of research pursues trainable chemistry in well-mixed reactors, where a reaction network plays the role of a neural network with reaction rates as trainable weights.
Neural-network behaviour, universal computation, and end-to-end training have been demonstrated \cite{hjelmfelt1991chemical, magnasco1997chemical, dack2024rncrn, cherry2018scaling}, but these systems are well-mixed: all information resides in scalar concentrations, with no spatial extent.
They sacrifice exactly the spatial continuity and parallelism that structured excitable media possess natively.
The present work keeps the spatial structure and asks how it can be abstracted and predicted.

\subsection{From Isolated Gates to Composed Computing}
\label{subsec:lit_gates}

The excitable-media computing literature is rich in isolated gates but poor in composition rules.
T\'{o}th and Showalter built a chemical XOR gate in which colliding pulses annihilate in a membrane-separated reaction--diffusion medium \cite{toth1995logic}.
Steinbock \etal used a geometrically structured BZ layer to implement logic gates and a chemical diode \cite{steinbock1996chemical}.
Adamatzky explored collision-based computing in excitable media, showing how streams of wave fragments can be routed and combined to perform computation \cite{adamatzky2002collision, adamatzky2002xor}.
G\'{o}recki \etal designed specific junction geometries to obtain logic gates and information-processing devices \cite{gorecki2009information}.

What is missing is a systematic way to compose these local gates into a predictive network model.
The present thesis addresses exactly this gap: it measures local transfer functions and composes them in a fast graph-level simulator, so that candidate networks can be explored before committing to expensive continuum simulations.

\section{The Belousov--Zhabotinsky Reaction and the Oregonator Model}
\label{sec:lit_bz}

The Belousov--Zhabotinsky (BZ) reaction, the metal-ion-catalysed oxidation of an organic substrate by bromate in acidic solution, is the canonical experimental excitable chemical medium.
It oscillates spontaneously under appropriate conditions and, in thin quasi-two-dimensional layers, supports propagating oxidation waves: expanding target patterns, rotating spirals, and interacting fronts.
Its phenomenology is well characterised, its ingredients are inexpensive and handled at room temperature, and compact models reproduce its behaviour faithfully.
It is therefore the standard testbed for chemical-wave computing.

The kinetic essence is captured by the Oregonator, Field and Noyes's three-variable reduction of the detailed Field--K\"or\"os--Noyes mechanism \cite{field1974oregonator}, and by Tyson and Fife's further reduction to two variables: a fast activator and a slow recovery variable \cite{tyson1980target}.
The singular-perturbation structure of the two-variable form underlies the classic understanding of pulse propagation in excitable media \cite{tyson1988singular}.
It is the model used throughout this thesis, coupled to spatial diffusion of the activator.

A variant of particular importance for computing is the photosensitive BZ reaction, in which the ruthenium-catalysed system is inhibited by illumination.
A spatial light field then acts as a writable mask: illuminated regions slow or block wave propagation, dark regions transmit it.
Sakurai \etal exploited this to route chemical waves through optically defined channels and junctions in real time \cite{sakurai2002design}.
In the language of this thesis, the light field $\phi(x,y)$ turns the substrate into a rewritable device.

The computing-relevant phenomenology of excitable media follows from a small set of robust behaviours.
A supra-threshold perturbation triggers a pulse travelling at a characteristic speed; behind it lies a refractory zone that cannot be re-excited, setting a maximum signalling frequency.
Colliding pulses annihilate; pulses meeting a narrowing gap can be blocked \cite{courtemanche1993instabilities}.
Structure adds more: asymmetric junctions act as chemical diodes \cite{agladze1996chemical}, T-shaped junctions as band-pass filters of pulse frequency \cite{gorecka2003tshaped}, and anisotropy creates preferred propagation directions \cite{steinbock1995anisotropy}.
Propagation, refractoriness, annihilation, block, rectification, filtering, and anisotropic routing are the native instruction set this work characterises.

\subsection{Pore-Network Models of Transport}
\label{subsec:lit_pore_networks}

The graph abstraction used in this thesis is borrowed from petroleum engineering and groundwater flow, where pore-network models predict multiphase flow and transport by assigning conductances to throats and solving for pressure and saturation on the graph \cite{xiong2016review, blunt2017multiphase}.
In those models a pore is a node, a throat is an edge, and the transport law is usually linear or weakly nonlinear.
The emphasis is on effective permeability, relative permeability curves, and dispersion, not on excitable pulses.

Pore-network models have also been applied to reactive transport, where species diffuse and react inside a network of pores \cite{meakin2009modeling}.
Again, the focus is on effective reaction rates rather than on threshold-activated pulse traffic.
Excitable media on networks have been studied by coupling discrete excitable elements according to a graph topology \cite{kinouchi2006optimal}, but the local element is typically a cellular-automaton node with a hand-chosen rule, not a calibrated continuous reaction--diffusion segment.

The present work differs in two ways.
First, the local element is a continuous reaction--diffusion segment whose transfer function is extracted from PDE simulations.
Second, the graph parameters are not chosen by hand but measured, so the model is calibrated rather than assumed.

\section{Synthesis and Research Gap}
\label{sec:lit_synthesis}

The literature reviewed above establishes five points.
First, physical neural networks are a viable alternative to digital inference, with trained demonstrations in optics, mechanics, and wave physics \cite{lin2018alloptical, stern2021physical, wright2022deep, hughes2019wave}.
Second, chemical and fluid media have a long computing pedigree, from liquid computers to bubble logic and chemical reservoirs \cite{adamatzky2019liquid, prakash2007bubble, baltussen2024chemical}.
Third, excitable media support wave-collision logic, diodes, filters, and routers, both in simulation and in BZ experiments \cite{toth1995logic, steinbock1996chemical, adamatzky2002xor, agladze1996chemical, gorecka2003tshaped}.
Fourth, well-mixed chemistry is trainable in the neural-network sense, but only by abandoning the spatial continuity that makes physical substrates attractive \cite{dack2024rncrn, cherry2018scaling}.
Fifth, molecular communication research treats reaction--diffusion transport as an information channel \cite{farsad2016comprehensive, guo2021molecular}, with non-coherent detectors and polynomial-chaos methods that quantify uncertainty \cite{lin2023diversity, abbaszadeh2019uncertainty}.

What is missing is a systematic way to move from the micro-scale physics inside a porous reactor to a predictive model of network-level pulse traffic.
Existing excitable-media work demonstrates isolated gates in ad hoc geometries; trainable-CRN work achieves learning but in a spaceless setting; and the physical-neural-network framework has not been carried into excitable chemical media arranged as a porous network.
Nobody has treated a structured BZ-type substrate as a set of local pore-scale operators, measured those operators in controlled two-dimensional simulations, and composed them into a fast graph-level predictor of pulse traffic.
This thesis addresses that gap.

\section{Research Aim}
\label{sec:aim}

The aim of this project is to characterise the local pulse-traffic rules inside a structured reaction--diffusion excitable medium and to show that they can be composed into a predictive pore-network model of a larger porous reactor.
Rather than asking whether such a medium can compute, the work maps out what it does natively at the pore scale, and how those local rules combine at the network scale.

\section{Research Objectives}
\label{sec:objectives}

To achieve this aim, the following objectives were set:

\begin{enumerate}
    \item Formulate the two-variable reaction--diffusion model for a two-dimensional domain, extended with structured boundaries, a spatial light-suppression field $\phi(x,y)$, and a spatially varying anisotropic diffusion tensor as substrate knobs.
    \item Implement an explicit finite-difference solver for this model and verify it against known BZ phenomenology and manufactured solutions.
    \item Design a black-box characterisation protocol in which the substrate is stimulated with pulses at input ports and its response is recorded at probe locations.
    \item Measure the native pore-scale transfer functions: channel pulse transmission, refractory-limited frequency response, two-input inhibition at a T-junction, and anisotropic routing.
    \item Build an event-driven graph simulator from the measured operators and validate it against full PDE solutions for simple networks.
    \item Identify the conditions under which the micro-scale operators transfer to three-dimensional pore geometries, and the limitations of the approach.
\end{enumerate}

\section{Scope and Delimitations}
\label{sec:scope}

In scope:
\begin{itemize}
    \item Two-dimensional reaction--diffusion dynamics in the Oregonator model, the chemically grounded reduction of the BZ mechanism.
    \item Explicit finite-difference time integration on a regular Cartesian grid.
    \item Three substrate parameterisations: wall geometry, the light-suppression field $\phi(x,y)$, and an anisotropic diffusion tensor field.
    \item Pulse-based input/output: localised activator perturbations as inputs, time-series probe recordings as outputs.
    \item A graph simulator built from measured local operators and validated against PDE solutions.
    \item Boundary three-dimensional tests that probe the transferability of the micro-scale operators.
\end{itemize}

Out of scope:
\begin{itemize}
    \item Full three-dimensional porous-medium simulations at the scale of a laboratory reactor.
    \item Experimental realisation or wet-lab validation.
    \item End-to-end training of a pore-network for a large computational task.
    \item Detailed comparison with acoustic, thermal, or well-mixed chemical alternatives; the starting point is the BZ reaction--diffusion medium.
\end{itemize}

The thesis is organised as follows.
\refChap{chap:porous} introduces the porous-reactor picture and the pore-network graph abstraction.
\refChap{chap:methodology} presents the Oregonator model, the finite-difference solver, and the graph-simulator construction.
\refChap{chap:results} reports the measured pore-scale operators and the graph-level validations.
\refChap{chap:conclusions} summarises the contributions, limitations, and next steps.
